% !TeX root = ../../main.tex

Abbildung~\ref{fig:experiment_flow} veranschaulicht den grundlegenden Ablauf eines Experiments.
Nach dem Start der Anwendung werden zunächst die für das jeweilige Experiment definierten Parameter geladen
und die entsprechenden Versuchsbedingungen festgelegt.
Dazu gehören unter anderem die Anzahl der durchzuführenden Transaktionen sowie die für den Versuch relevanten
Konfigurationsparameter.
Anschließend wird das Experiment gestartet und die einzelnen Transaktionen entsprechend der vorgegebenen Konfiguration
ausgeführt.
Für jede Transaktion wird zunächst eine SQL-Anfrage an die Datenbank gesendet und ausgeführt.
Nach der Ausführung wird das Ergebnis der Transaktion ausgewertet.
Dabei wird insbesondere geprüft, ob die Transaktion erfolgreich abgeschlossen wurde oder ob ein Fehler aufgetreten ist,
der einen erneuten Ausführungsversuch erforderlich macht.
Tritt ein entsprechender Fehler auf und steht gemäß der definierten Retry-Strategie noch ein weiterer Versuch zur
Verfügung, wird die Transaktion erneut ausgeführt.
Dieser Ablauf kann sich abhängig vom Ergebnis und der maximal zulässigen Anzahl an Wiederholungsversuchen
mehrfach wiederholen.
Sobald die Transaktion erfolgreich abgeschlossen wurde oder keine weiteren Versuche mehr möglich sind, wird mit der
nächsten Transaktion fortgefahren.
Nach Abschluss aller für das Experiment vorgesehenen Transaktionen werden die während der Ausführung erfassten
Ergebnisse und Messwerte gespeichert.
Hierzu gehören hauptsächlich die Informationen über erfolgreiche und fehlgeschlagene Transaktionen sowie die Anzahl
der durchgeführten Wiederholungsversuche.
Auf Grundlage dieser aufgezeichneten Daten werden anschließend die für die Auswertung relevanten Kenngrößen berechnet.
Dadurch kann das Verhalten des untersuchten Systems unter den jeweiligen Versuchsbedingungen nachvollziehbar
erfasst und anschließend mit den Ergebnissen anderer Experimente verglichen werden.

% !TeX root = ../main.tex

\begin{figure}[ht]
	\centering
	\begin{tikzpicture}
		\node[
			fill=TerminalBackground,
			draw=TerminalFrame,
			rounded corners,
			inner sep=10pt
		] {
			\begin{tikzpicture}[
				scale=0.8,
				transform shape,
				node distance=1cm,
				every node/.style={font=\small},
				startstop/.style={
					rectangle,
					rounded corners,
					draw,
					minimum width=3cm,
					minimum height=1cm,
					align=center
				},
				process/.style={
					rectangle,
					draw,
					minimum width=3cm,
					minimum height=1cm,
					align=center
				},
				decision/.style={
					diamond,
					draw,
					align=center,
					aspect=2
				},
				arrow/.style={
					->,
					>=stealth
				}
			]

			\node (start) [startstop] {Anwendung starten};
			\node (error) [process, below=of start] {Namensgebende Funktion\\ startet Experimente};
			\node (experiment) [process, below=of error] {Experiment startet Retry\\ anhand der definierten Parameter};
			\node (run) [process, below=of experiment] {SQL ausführen};
			\node (sOrf) [decision, below=of run] {success or failure};
			\node (possible) [decision, below right=1.5cm of sOrf] {Retry möglich?};
			\node (retry) [process, below right=2cm of possible] {Retry ausführen};
			\node (log) [process, below=3cm of sOrf] {Log result};
			\node (calc) [process, below=of log] {Kenngrößen berechnen};
			\node (result) [process, below=of calc] {Ergebnisse auswerten};
			\node (end) [startstop, below=of result] {Experiment beenden};

			\draw [arrow, shorten <=5pt, shorten >=5pt] (start) -- (error);
			\draw [arrow, shorten <=5pt, shorten >=5pt] (error) -- (experiment);
			\draw [arrow, shorten <=5pt, shorten >=5pt] (experiment) -- (run);
			\draw [arrow, shorten <=5pt, shorten >=5pt] (run) -- (sOrf);
			\draw [arrow, shorten <=5pt, shorten >=5pt] (sOrf) -- node[midway, above right] {failed} (possible);
			\draw [arrow, shorten <=5pt, shorten >=5pt] (sOrf) -- node[midway, above left] {succeeded} (log);
			\draw [arrow, shorten <=5pt, shorten >=5pt] (possible) -- node[below left] {possible} (retry);
			\draw [arrow, shorten <=5pt, shorten >=5pt] (possible) -- node[below right] {not possible} (log);
			\draw [arrow, shorten <=5pt, shorten >=5pt] (retry) |- (sOrf);
			\draw [arrow, shorten <=5pt, shorten >=5pt] (log) -- (calc);
			\draw [arrow, shorten <=5pt, shorten >=5pt] (calc) -- (result);
			\draw [arrow, shorten <=5pt, shorten >=5pt] (result) -- (end);

			\end{tikzpicture}};
	\end{tikzpicture}
	
	\caption{Ablauf eines Experiments}
	\label{fig:experiment_flow}
\end{figure}